% AUTO-GENERATED by Github-Branch/tools/build_unified_source.ps1.
% Do not edit this file directly: update the transformation manifest instead.
\chapter{Econometrics and Time-Series Foundations}
\chapterstatus{Edited reuse from Team 2. General theory is retained; application-specific claims require the unified validation protocol.}
\sourceassetnote
\section{The Simple Linear Regression Model and Ordinary Least Squares}
% ===============================

Building on the foundational concepts, we now delve into the Simple Linear Regression (SLR) model, the most basic framework for analyzing the relationship between two variables. Understanding the SLR model provides essential intuition for more complex analyses.
\subsection{Defining the Population Model and Core Assumptions}

The SLR model aims to explain a dependent variable $y$ using a single independent variable $x$. It relies on several key assumptions about the underlying population relationship and the data generation process. We label these assumptions with the prefix SLR to distinguish them from the more general multiple regression assumptions introduced later.

\textbf{Assumption SLR.1 (Linear in Parameters):} The model posits a linear relationship in the population:
$$y = \beta_0 + \beta_1 x + u$$

Here, $\beta_0$ is the intercept parameter and $\beta_1$ is the slope parameter, representing the change in $y$ for a one-unit increase in $x$, holding all other factors constant. The term $u$ is the error term, capturing all factors affecting $y$ other than $x$.

\textbf{Assumption SLR.2 (Random Sampling):} We work with a dataset containing $n$ pairs of observations $\{(x_i, y_i): i = 1, \ldots, n\}$, randomly and independently drawn (i.i.d.) from the population that follows the linear relationship. Random sampling guarantees that our dataset is representative of the population, allowing us to use the sample data to estimate the true population parameters.

\textbf{Assumption SLR.3 (Sample Variation in the Explanatory Variable):} The sample outcomes on $x$ are not all the same value. This minimal requirement ensures that we have enough variation in the independent variable to estimate its relationship with $y$. Without variation in $x$, the OLS slope estimator $\hat{\beta}_1$ is undefined.

\textbf{Assumption SLR.4 (Zero Conditional Mean):} This assumption is crucial for interpreting the estimated coefficient $\beta_1$ as a causal effect. It states that the expected value of the error term $u$, given any value of the explanatory variable $x$, is zero:
$$E(u \mid x) = 0$$

In other words, the unobserved factors captured by $u$ should not be systematically related to $x$. Their average effect must be zero no matter what value $x$ takes. If instead $E(u \mid x)$ changes with $x$, it means that $x$ is somehow linked to factors not included in the model, and OLS cannot separate the true effect of $x$ on $y$ from the influence of those missing factors.

Writing down the model $y = \beta_0 + \beta_1 x + u$ is only the starting point. What really matters is whether the explanatory variable $x$ is independent of the error term $u$. If they are not related on average, OLS works correctly and gives an unbiased estimate of how $x$ affects $y$. However, if this assumption fails---for example, because we omitted some important variable that is correlated with $x$---then our estimate of $\beta_1$ will be misleading. This situation is known as \textbf{omitted variable bias}, and it is one of the main threats to causal interpretation in regression models.

Consider the simple gold pricing model:
$$\text{gold price} = \beta_0 + \beta_1 \text{real rates} + u$$

where $u$ includes unobserved factors such as inflation expectations, geopolitical risks, and central bank policies. The assumption $E(u \mid \text{real rates}) = 0$ implies that, on average, these unobserved factors do not vary with real interest rate levels. However, when inflation expectations rise, both real interest rates (set higher by central banks to combat inflation) and gold prices (sought as an inflation hedge) increase together. In that case, inflation expectations and real rates are positively correlated, meaning $E(u \mid \text{real rates})$ depends on real rates. As a result, $\hat{\beta}_1$ would misattribute some of the correlation between rates and gold to the direct causal effect---a clear case of omitted variable bias.

When SLR.4 holds (and $E(u) = 0$, which can be assumed without loss of generality if an intercept $\beta_0$ is included), we can define the \textbf{Population Regression Function (PRF)}:
$$E(y \mid x) = \beta_0 + \beta_1 x$$

The PRF shows how the average value of $y$ changes linearly with $x$ in the population. It represents the systematic part of $y$, while the error term $u$ captures the unsystematic part.
\subsection{Estimating the Parameters: Ordinary Least Squares (OLS)}

The parameters $\beta_0$ and $\beta_1$ describe the population relationship but are unknown. To estimate them, we collect data via random sampling (SLR.2), obtaining a sample $\{(x_i, y_i): i = 1, \ldots, n\}$.

The most widely used method is \emph{Ordinary Least Squares (OLS)}. OLS works by finding the line through the scatterplot of the sample data points that minimizes the sum of the squared vertical distances from each point to the line.

For any candidate line $\hat{y} = b_0 + b_1 x$, the predicted value for observation $i$ is $\hat{y}_i = b_0 + b_1 x_i$ (the \emph{fitted value}), and the prediction error is $\hat{u}_i = y_i - \hat{y}_i$ (the \emph{residual}). OLS selects the specific estimates $\hat{\beta}_0$ and $\hat{\beta}_1$ that minimize the \emph{Sum of Squared Residuals (SSR)}:
$$\min_{b_0, b_1} \sum_{i=1}^{n} (y_i - b_0 - b_1 x_i)^2$$

Using calculus (or the method of moments based on $E(u) = 0$ and $E(xu) = 0$), the OLS estimators are:
$$\hat{\beta}_1 = \frac{\sum_{i=1}^{n}(x_i - \bar{x})(y_i - \bar{y})}{\sum_{i=1}^{n}(x_i - \bar{x})^2}, \qquad \hat{\beta}_0 = \bar{y} - \hat{\beta}_1 \bar{x}$$

where $\bar{x}$ and $\bar{y}$ are the sample means. The slope estimate $\hat{\beta}_1$ is the sample covariance between $x$ and $y$ divided by the sample variance of $x$.

The resulting estimated line, $\hat{y} = \hat{\beta}_0 + \hat{\beta}_1 x$, is called the OLS regression line or the \emph{Sample Regression Function (SRF)}, which serves as our estimate of the unknown PRF.
\subsection{Algebraic Properties of OLS and Goodness-of-Fit}

The OLS estimation process guarantees certain algebraic properties that hold for any sample of data:
\begin{enumerate}
  \item The sum of the OLS residuals is exactly zero: $\sum_{i=1}^{n} \hat{u}_i = 0$.
  \item The sample covariance between $x_i$ and the OLS residuals $\hat{u}_i$ is zero: $\sum_{i=1}^{n} x_i \hat{u}_i = 0$.
  \item The point $(\bar{x}, \bar{y})$ always lies exactly on the OLS regression line.
\end{enumerate}

These properties are direct consequences of the minimization procedure. OLS separates each $y_i$ into what the model explains ($\hat{y}_i$) and what remains unexplained ($\hat{u}_i$). These two parts are uncorrelated within the sample.

This decomposition allows us to measure how well the regression line fits the data. The \emph{Total Sum of Squares (SST)} represents the total variation of $y$:
$$SST = \sum_{i=1}^n (y_i - \bar{y})^2$$

It decomposes into the \emph{Explained Sum of Squares (SSE)} and the \emph{Sum of Squared Residuals (SSR)}:
$$SSE = \sum_{i=1}^n (\hat{y}_i - \bar{y})^2, \qquad SSR = \sum_{i=1}^n \hat{u}_i^2, \qquad SST = SSE + SSR$$

The $R^2$, or coefficient of determination, measures the proportion of total variation in $y$ explained by the regression:
$$R^2 = \frac{SSE}{SST} = 1 - \frac{SSR}{SST}$$

$R^2$ is always between 0 and 1. While it provides a summary measure of fit, it is important not to overemphasize $R^2$: a model can have a low $R^2$ but still provide a reliable estimate of the ceteris paribus effect $\beta_1$, especially if the goal is causal inference rather than pure prediction.

\section[Statistical Properties: Unbiasedness, Variance, and the Gauss--Markov Theorem]{Statistical Properties: Unbiasedness, Variance, and\\ the Gauss--Markov Theorem}

We now transition from describing the properties of OLS estimates within a single sample to examining the \textbf{statistical properties} of the OLS estimators across repeated random sampling from the population. Rather than treating $\hat{\beta}_0$ and $\hat{\beta}_1$ as fixed numbers computed from one dataset, we view them as \textbf{random variables} whose realizations depend on which particular sample we happen to draw. The distribution of these estimates across all possible samples is called the \textbf{sampling distribution}, and understanding its properties is essential for making reliable statistical inferences.

A cornerstone property for any estimator is \emph{unbiasedness}. Under the \textbf{first four SLR assumptions} (SLR.1 through SLR.4), the OLS estimators possess this crucial property:

\begin{theorem}[Unbiasedness of OLS Estimators]
$$E(\hat{\beta}_0) = \beta_0 \quad \text{and} \quad E(\hat{\beta}_1) = \beta_1$$
\end{theorem}

This means that if many random samples were drawn and the OLS estimates computed each time, the average of these estimates would equal the true population values. Unbiasedness is a statement about the \textbf{procedure} over repeated samples, not about any single estimate. However, this crucial property relies on \textbf{SLR.4}; if $E(u \mid x) \neq 0$, the OLS estimators will generally be \textbf{biased}.

While unbiasedness tells us that OLS targets the true parameter on average, it says nothing about the \textbf{precision} of the sampling distribution. Precision is measured by the \textbf{sampling variance} of the estimators. To derive tractable formulas for these variances, we introduce:

\textbf{Assumption SLR.5 (Homoskedasticity):}
$$\operatorname{Var}(u \mid x) = \sigma^2$$

This states that the conditional variance of the error term is \emph{constant} across all values of $x$. If the variance changes with $x$---a situation called \textbf{heteroskedasticity}---the data violates this assumption. Heteroskedasticity means that some observations are more noisy than others: for example, in financial data, volatility often varies with market conditions or asset prices, so the spread of returns around their conditional mean is not constant over time.

Under the complete set of \textbf{Gauss-Markov Assumptions} (SLR.1 through SLR.5), the variances of the OLS estimators take on particularly simple forms:
$$\operatorname{Var}(\hat{\beta}_1) = \frac{\sigma^2}{SST_x}, \qquad \operatorname{Var}(\hat{\beta}_0) = \sigma^2\left(\frac{1}{n} + \frac{\bar{x}^2}{SST_x}\right)$$

where $SST_x = \sum_{i=1}^{n}(x_i - \bar{x})^2$. These formulas reveal that the precision of $\hat{\beta}_1$ improves (variance decreases) when there is more variation in $x$ and when the underlying error variance $\sigma^2$ is smaller.

A fundamental result of classical econometrics is:

\begin{theorem}[Gauss-Markov Theorem]
Under assumptions SLR.1 through SLR.5, the OLS estimators $\hat{\beta}_0$ and $\hat{\beta}_1$ are the \textbf{Best Linear Unbiased Estimators (BLUE)}. That is, among all linear unbiased estimators, OLS has the smallest variance.
\end{theorem}

Being ``Best'' means minimum variance; being ``Linear Unbiased'' means no systematic bias. This theorem provides the formal justification for using OLS: under the stated assumptions, there is no better linear estimator.

Since $\sigma^2$ is unknown, it must be estimated from the data. An unbiased estimator is:
$$\hat{\sigma}^2 = \frac{SSR}{n - 2}$$

The denominator is $n - 2$ (not $n$) because we lose two \emph{degrees of freedom} by estimating two parameters ($\beta_0$ and $\beta_1$). Using $n$ instead would systematically underestimate $\sigma^2$; dividing by $n - 2$ corrects for this and yields an unbiased, consistent estimator.

The positive square root, $\hat{\sigma}$, is called the \textbf{Standard Error of the Regression (SER)} and estimates the standard deviation of the error term $u$. The standard errors of the OLS estimates are:
$$se(\hat{\beta}_1) = \frac{\hat{\sigma}}{\sqrt{SST_x}}, \qquad se(\hat{\beta}_0) = \hat{\sigma}\sqrt{\frac{1}{n} + \frac{\bar{x}^2}{SST_x}}$$

These standard errors quantify the uncertainty surrounding each estimated coefficient and are fundamental for constructing confidence intervals and test statistics.

% ===============================

% ===============================
\section{Multiple Regression, Inference, and Large-Sample Methods}
% ===============================

This chapter brings together the logic of multiple regression, finite-sample inference, and large-sample approximation in a single progression. The objective is to keep the transition from interpretation to testing and asymptotic justification continuous, so that the reader can move from economic meaning to statistical inference without artificial breaks.

In financial markets, variables such as real interest rates, the dollar value, and equity indices tend to move together. Reading the effect of a single variable in a bivariate model risks confounding co-movements with causality. Multiple regression introduces a framework in which the effect of each predictor is evaluated \textbf{holding other conditions constant}, allowing the isolation of information components that would otherwise remain intertwined.
\subsection{From Motivation to Method: The Case for Multiple Predictors}

The need to consider multiple regressors simultaneously arises from the necessity to separate overlapping signals. Suppose we estimate how gold returns vary with real interest rates and the dollar value. If real rates and the exchange rate move together, the estimated effect of each will be difficult to isolate: the rate coefficient will capture both the contribution of the rate itself and part of the uncontrolled exchange rate effect. Adding a third variable to control for this overlap improves interpretive precision.

Formally, the multiple linear regression (MLR) model is:
\begin{equation}
y_t = \beta_0 + \beta_1 x_{1t} + \beta_2 x_{2t} + \cdots + \beta_k x_{kt} + u_t
\end{equation}

where $y_t$ is the dependent variable, $\mathbf{x}_t = (x_{1t}, \ldots, x_{kt})'$ is the vector of $k$ regressors, $\beta_0, \ldots, \beta_k$ are the parameters of interest, and $u_t$ is the error term. The crucial condition for correct interpretation is \textbf{exogeneity}: $E(u_t \mid \mathbf{x}_t) = 0$, meaning the absence of systematic dependence between unobserved factors and the regressors. Under this assumption, the coefficients acquire causal-interpretive significance and are not merely correlational.
\subsection{Matrix Form of the Multiple Regression Model}

For compact notation and analytical tractability, the MLR model is most elegantly expressed in matrix form. Define the $n \times (k+1)$ design matrix $\mathbf{X}$, the $n \times 1$ response vector $\mathbf{y}$, the $(k+1) \times 1$ parameter vector $\boldsymbol{\beta}$, and the $n \times 1$ error vector $\mathbf{u}$:
$$\mathbf{X} = \begin{pmatrix} 1 & x_{11} & \cdots & x_{1k} \\ 1 & x_{21} & \cdots & x_{2k} \\ \vdots & \vdots & \ddots & \vdots \\ 1 & x_{n1} & \cdots & x_{nk} \end{pmatrix}, \quad \mathbf{y} = \begin{pmatrix} y_1 \\ y_2 \\ \vdots \\ y_n \end{pmatrix}, \quad \boldsymbol{\beta} = \begin{pmatrix} \beta_0 \\ \beta_1 \\ \vdots \\ \beta_k \end{pmatrix}, \quad \mathbf{u} = \begin{pmatrix} u_1 \\ u_2 \\ \vdots \\ u_n \end{pmatrix}$$

The entire system of $n$ equations is:
$$\mathbf{y} = \mathbf{X}\boldsymbol{\beta} + \mathbf{u}$$

The OLS estimator minimizes $\mathbf{u}'\mathbf{u} = (\mathbf{y} - \mathbf{X}\boldsymbol{\beta})'(\mathbf{y} - \mathbf{X}\boldsymbol{\beta})$, yielding the closed-form solution:
$$\hat{\boldsymbol{\beta}} = (\mathbf{X}'\mathbf{X})^{-1}\mathbf{X}'\mathbf{y}$$

provided that $\mathbf{X}'\mathbf{X}$ is invertible (which requires no perfect multicollinearity among the regressors). The variance-covariance matrix of $\hat{\boldsymbol{\beta}}$ under homoskedasticity is:
$$\operatorname{Var}(\hat{\boldsymbol{\beta}} \mid \mathbf{X}) = \sigma^2 (\mathbf{X}'\mathbf{X})^{-1}$$
\subsection{Partial Effects and Economic Interpretation}

The coefficient $\beta_j$ represents \textbf{the average change in $y_t$ associated with a one-unit change in $x_{jt}$, holding all other regressors constant}. In our gold price example, the real interest rate coefficient measures how the price moves as the rate varies, provided the dollar value and other factors remain fixed.

When data undergo logarithmic transformations, the meaning of coefficients changes in form: $\beta_j$ becomes an elasticity if both the regressor and the dependent variable are in logs, or a semi-elasticity if only one is. A coefficient of 0.03 on a log-regressor in a log-linear model means that a 1\% increase in the regressor is associated with a 0.03\% increase in the dependent variable.

An important distinction must be kept in mind: a coefficient may be \textbf{statistically significant} yet \textbf{economically negligible}. Statistical significance depends on both the magnitude of the coefficient and the precision of estimation, while economic significance depends solely on the magnitude relative to the scale of the variables. With large samples, even trivially small effects can achieve statistical significance. One must always evaluate both dimensions.
\subsection{The Frisch-Waugh-Lovell Theorem}

The \textbf{Frisch-Waugh-Lovell (FWL) theorem} provides an elegant operational translation of partial effects. The partial effect of $x_j$ is obtained by purging both the dependent variable and the regressor of interest from all other information contained in the remaining predictors:
\begin{enumerate}
  \item Regress $y_t$ on $\mathbf{x}_{-j,t}$ (all regressors except $x_{jt}$) and save the residuals $\tilde{y}_t$.
  \item Regress $x_{jt}$ on $\mathbf{x}_{-j,t}$ and save the residuals $\tilde{x}_{jt}$.
  \item Regress $\tilde{y}_t$ on $\tilde{x}_{jt}$ without an intercept; the estimated coefficient is exactly $\hat{\beta}_j$ from the full multiple regression.
\end{enumerate}

This theorem shows that a coefficient in multiple regression can be viewed as the result of a univariate regression between the residual components of each variable after removing the influence of all other predictors. This clarifies why the sign and significance of $\hat{\beta}_j$ can change when regressors are added or removed: the unique component of $x_j$ unexplained by other predictors changes.
\subsection{Multicollinearity: When Informational Overlap Becomes Problematic}

Strong co-movements between regressors do not create systematic distortion (bias) in coefficients but instead amplify their variance. This phenomenon, known as \textbf{multicollinearity}, manifests through wide standard errors, weak tests of statistical significance, and instability: small changes in the sample can cause marked variations in estimated coefficients.

Practical diagnosis of multicollinearity combines three complementary indicators:
\begin{itemize}
  \item \textbf{Variance Inflation Factor (VIF)}: for each regressor $x_j$, the VIF measures how much the variance of $\hat{\beta}_j$ is inflated by the correlation of $x_j$ with other regressors:
  $$\text{VIF}_j = \frac{1}{1 - R_j^2}$$
  where $R_j^2$ is the $R^2$ from regressing $x_j$ on all other regressors. A high VIF (commonly, above 5--10) signals informational redundancy.
  \item \textbf{Pairwise correlations}: direct inspection of correlations between pairs of regressors; values near $\pm 1$ indicate tight co-movement.
  \item \textbf{Condition number} of the $\mathbf{X}'\mathbf{X}$ matrix: a synthetic measure of ill-conditioning; large values (beyond 30--100) indicate proximity to exact collinearity.
\end{itemize}
\subsection{Specification Choices: Parsimony, Overfitting, and Model Selection}

Managing multicollinearity rests on two complementary levers.

\textbf{First, parsimonious specification}: when two variables tell nearly the same economic story, one selects the better-measured, more interpretable, or more relevant variable and omits the other. This choice must be theoretically motivated, not arbitrary.

\textbf{Second, variable transformations when economically justified:} logarithms, deflating nominal quantities, or other transformations can align variables more meaningfully. Such transformations should be motivated by economic theory, not merely as a technical fix.

A related concern is \textbf{overfitting}: adding regressors to improve in-sample fit without genuine predictive power. An overfitted model performs well on the estimation sample but poorly out-of-sample. To balance fit against parsimony, information criteria provide principled model selection:
$$\text{AIC}(k) = -2\log(\hat{L}) + 2k, \qquad \text{BIC}(k) = -2\log(\hat{L}) + k\log(n)$$

where $\hat{L}$ is the maximized likelihood, $k$ is the number of parameters, and $n$ is the sample size. The BIC penalizes complexity more heavily for large $n$. The preferred model minimizes the chosen criterion.
\subsection{Uncertainty and Robust Inference in Financial Data}

Economic and financial time series frequently depart from the ideal setting of independent and homoskedastic errors. Heteroskedasticity (non-constant variance) and autocorrelation (nearby errors are correlated in time) are ubiquitous. Ignoring these features leads to unreliable uncertainty measures.

To correct this problem, we resort to \textbf{robust standard errors}:
\begin{itemize}
    \item \textbf{Heteroskedasticity-consistent (HC) standard errors}:\\
    When the error variance varies across observations, the usual OLS standard-error formula is no longer valid. White's heteroskedasticity-consistent correction yields the so-called ``sandwich'' estimator:
  $$\widehat{\operatorname{Var}}(\hat{\boldsymbol{\beta}}) = (\mathbf{X}'\mathbf{X})^{-1} \left(\sum_{i=1}^n \hat{u}_i^2\, \mathbf{x}_i\mathbf{x}_i'\right) (\mathbf{X}'\mathbf{X})^{-1}$$
  This estimator is valid under heteroskedasticity of any form, as long as the model is correctly specified.
  \item \textbf{Heteroskedasticity and Autocorrelation Consistent (HAC) standard errors}: In time series data, errors often exhibit autocorrelation. The Newey-West procedure adjusts standard errors to account for both heteroskedasticity and temporal dependence, with the lag window chosen based on data frequency and economic reasoning about persistence.
\end{itemize}

Once robust standard errors are computed, it is preferable to report \textbf{confidence intervals} rather than restricting oneself to binary p-values. A wide interval signals high uncertainty; a narrow interval communicates precision. Understanding this dynamic is fundamental to avoiding misinterpretations, particularly when high $R^2$ coexists with imprecise coefficients due to collinearity.

% ===============================

In the previous section, we defined the MLR model and established the Gauss--Markov assumptions. Under these assumptions, OLS estimators are BLUE. We derived formulas for the sampling variances of the OLS estimators. However, unbiasedness and variance formulas alone are not sufficient for hypothesis testing. To perform statistical inference, we need to know the \textbf{full sampling distribution} of the OLS estimators.
\subsection{The Normality Assumption and the Classical Linear Model}

To make sampling distributions tractable, we add one final assumption.

\textbf{Assumption MLR.6 (Normality):} The population error $u$ is independent of the explanatory variables $x_1, \ldots, x_k$ and is normally distributed:
$$u \sim \text{Normal}(0, \sigma^2)$$

This assumption is much stronger than the others. It implies both zero conditional mean and homoskedasticity. The set of assumptions MLR.1 through MLR.6 constitutes the \textbf{Classical Linear Model (CLM)} assumptions. Under the CLM, the model can be written as:
$$y \mid \mathbf{x} \sim \text{Normal}(\beta_0 + \beta_1 x_1 + \cdots + \beta_k x_k, \sigma^2)$$

It is important to understand the role of this assumption. The normality assumption is what delivers \textbf{exact inference}: if MLR.6 holds, then $t$-statistics follow exact $t$-distributions and $F$-statistics follow exact $F$-distributions in \emph{any} sample size. In contrast, without MLR.6, we can appeal to the \textbf{Central Limit Theorem} to show that these statistics follow the same distributions \emph{approximately} in \textbf{large samples} (asymptotic inference). The distinction matters in small samples: exact inference is valid for any $n$, while asymptotic inference is only valid when $n$ is large enough.

The normality of $u$ can be loosely justified by the CLT: if $u$ is the sum of many independent unobserved factors, the aggregate may approximate a normal distribution. However, in financial applications, asset returns often exhibit heavy tails and skewness, making the normality assumption questionable. This is why large-sample methods are so important in finance.

\begin{theorem}[Normal Sampling Distributions under CLM]
Under the CLM assumptions, the OLS estimators $\hat{\beta}_j$ are normally distributed conditional on the sample values of the explanatory variables:
$$\hat{\beta}_j \sim \text{Normal}\left(\beta_j, \operatorname{Var}(\hat{\beta}_j)\right)$$
\end{theorem}
\subsection{Testing Hypotheses about a Single Parameter: The \texorpdfstring{$t$}{t}-Test}

The normal sampling distribution requires knowing the true population standard deviation $\sigma(\hat{\beta}_j)$, which depends on the unknown $\sigma$. In practice, we replace $\sigma$ with its estimate $\hat{\sigma}$ (the SER). This replacement changes the distribution from normal to a $t$-distribution.

\begin{theorem}[$t$-Distribution under CLM]
Under the CLM assumptions:
$$\frac{\hat{\beta}_j - \beta_j}{se(\hat{\beta}_j)} \sim t_{n-k-1}$$
\end{theorem}

The degrees of freedom are $n - (k+1)$: we subtract $k+1$ because we have estimated $k+1$ parameters (the $k$ slopes plus the intercept). Each estimated parameter uses up one degree of freedom. The greater the number of estimated parameters relative to the sample size, the less information remains for inference, and the wider the $t$-distribution becomes.

To test the null hypothesis $H_0: \beta_j = 0$, we compute the $t$-statistic:
$$t_{\hat{\beta}_j} = \frac{\hat{\beta}_j}{se(\hat{\beta}_j)}$$

This measures how many estimated standard deviations $\hat{\beta}_j$ is away from zero. We compare this value to a critical value $c$ from the $t_{n-k-1}$ distribution:

\begin{itemize}
\item \textbf{Two-Sided Alternative} ($H_1: \beta_j \neq 0$): Reject $H_0$ if $|t_{\hat{\beta}_j}| > c$, where $c$ is the 97.5th percentile for a 5\% test.
\item \textbf{One-Sided Alternative} ($H_1: \beta_j > 0$): Reject $H_0$ if $t_{\hat{\beta}_j} > c$.
\item \textbf{One-Sided Alternative} ($H_1: \beta_j < 0$): Reject $H_0$ if $t_{\hat{\beta}_j} < -c$.
\end{itemize}

It is crucial to distinguish between \textbf{statistical significance} and \textbf{economic significance}. Statistical significance is determined by the size of the $t$-statistic, which depends on both the coefficient's magnitude and the standard error's smallness. Economic significance is determined by the \textbf{magnitude and sign of $\hat{\beta}_j$} relative to the scale of the variables. With large sample sizes, standard errors become very small, and even trivially small effects can achieve statistical significance. Therefore, after checking for statistical significance, one must always interpret the magnitude of the coefficient to assess its practical importance.
\subsection{Confidence Intervals}

Instead of just testing a hypothesis, we can construct a \textbf{confidence interval (CI)}, which provides a range of plausible values for the unknown population parameter $\beta_j$. A 95\% confidence interval is constructed as:
$$\hat{\beta}_j \pm c \cdot se(\hat{\beta}_j)$$

where $c$ is the 97.5th percentile in the $t_{n-k-1}$ distribution.

The correct interpretation is based on repeated sampling: if we were to draw many samples and construct a 95\% CI for each, \textbf{95\% of those intervals would contain the true population value $\beta_j$}. The CI does \emph{not} mean that the true parameter lies in this specific interval with 95\% probability---the true parameter is fixed; it is the interval that is random, varying from sample to sample.

Confidence intervals provide a direct link to two-sided hypothesis tests: we reject $H_0: \beta_j = a_j$ at the 5\% level if and only if $a_j$ is not contained within the 95\% CI. Reporting CIs is generally more informative than reporting only p-values, because they convey both the direction and precision of the estimated effect.
\subsection{Testing Linear Combinations of Parameters}

Sometimes we are interested in testing a hypothesis involving multiple parameters, such as $H_0: \beta_1 = \beta_2$. This can be rewritten as $H_0: \beta_1 - \beta_2 = 0$. To test this, we compute:
$$t = \frac{(\hat{\beta}_1 - \hat{\beta}_2)}{se(\hat{\beta}_1 - \hat{\beta}_2)}, \quad se(\hat{\beta}_1 - \hat{\beta}_2) = \sqrt{[se(\hat{\beta}_1)]^2 + [se(\hat{\beta}_2)]^2 - 2s_{12}}$$

where $s_{12}$ estimates $\operatorname{Cov}(\hat{\beta}_1, \hat{\beta}_2)$.

A simpler method is to \textbf{reparameterize the model}. Defining $\theta_1 = \beta_1 - \beta_2$ (so $\beta_1 = \theta_1 + \beta_2$) and substituting:
$$y = \beta_0 + \theta_1 x_1 + \beta_2(x_1 + x_2) + \ldots + u$$

By running this reparameterized regression, the OLS output for $x_1$ directly provides $\hat{\theta}_1$ and $se(\hat{\theta}_1)$, enabling a standard $t$-test.

\section[Testing Multiple Restrictions: F-Test and Wald Statistic]
{Testing Multiple Restrictions: \texorpdfstring{$F$}{F}-Test and Wald Statistic}

Often we need to test a \textbf{joint hypothesis} involving multiple restrictions simultaneously, such as $H_0: \beta_3 = 0, \beta_4 = 0, \beta_5 = 0$. Using individual $t$-tests for each parameter is not valid; variables that are individually insignificant (perhaps due to multicollinearity) can be highly significant as a group.

The \textbf{$F$-test} compares the fit of the \textbf{unrestricted model} (with all variables) and the \textbf{restricted model} (where $H_0$ is imposed):
$$F = \frac{(SSR_r - SSR_u)/q}{SSR_u/(n - k - 1)}$$

where $q$ is the number of restrictions (numerator degrees of freedom). The $F$-statistic measures whether the increase in SSR when moving from the unrestricted to the restricted model is large enough to justify rejecting $H_0$. Under $H_0$ and the CLM assumptions, $F \sim F_{q, n-k-1}$.

The probabilistic interpretation of the $F$-test is illuminating: the $p$-value is the probability of observing an $F$-statistic as large as the one computed, \emph{assuming the null hypothesis is true}. A small $p$-value means the data are unlikely under $H_0$. Equivalently, the numerator measures the average reduction in SSR per restriction: if restricting the model significantly worsens fit, we reject the null.

In the general case of $q$ linear restrictions $\mathbf{R}\boldsymbol{\beta} = \mathbf{r}$, the \textbf{Wald statistic} takes the matrix form:
$$W = (\mathbf{R}\hat{\boldsymbol{\beta}} - \mathbf{r})'\left[\mathbf{R}(\mathbf{X}'\mathbf{X})^{-1}\mathbf{R}'\hat{\sigma}^2\right]^{-1}(\mathbf{R}\hat{\boldsymbol{\beta}} - \mathbf{r})$$

Under $H_0$, $W/q \sim F_{q, n-k-1}$. The Wald statistic nests all standard $F$-tests as special cases.

% ===============================

Up to this point, our analysis has focused on finite sample properties of OLS, which rely on assumptions holding for any sample size $n$, and crucially, on the normality of the error term (MLR.6) for exact inference. In many empirical applications, particularly with financial data, error terms may not be normally distributed. In such cases, we rely on \textbf{asymptotic properties}: properties of estimators that hold as the sample size grows indefinitely.
\subsection{Consistency of the OLS Estimators}

\textbf{Consistency} is a minimal requirement for any useful estimator in large samples. An estimator $\hat{\beta}_j$ is consistent if, as $n \to \infty$, its distribution collapses to the single point $\beta_j$. This convergence in probability is written formally as:
$$\plim_{n \to \infty} \hat{\beta}_j = \beta_j$$

The formal foundation of consistency is the \textbf{Law of Large Numbers (LLN)}: for an i.i.d. sequence with finite mean, the sample average converges in probability to the population mean as $n \to \infty$. Applying the LLN to the OLS normal equations, we can show:

\begin{theorem}[Consistency of OLS]
Under the first four Gauss-Markov assumptions (MLR.1--MLR.4), the OLS estimator $\hat{\beta}_j$ is consistent for $\beta_j$:
$$\plim \hat{\beta}_j = \beta_j$$
\end{theorem}

While unbiasedness requires $E(u \mid x_1, \ldots, x_k) = 0$, consistency only requires $\operatorname{Cov}(x_j, u) = 0$ for all $j$---a weaker condition. However, if $x_j$ is correlated with $u$ (e.g., due to omitted variables), OLS is not only biased but also \textbf{inconsistent}: the bias does not vanish as $n$ grows.

For a simple regression model $y = \beta_0 + \beta_1 x_1 + u$, if $x_1$ is correlated with $u$:
$$\plim \hat{\beta}_1 = \beta_1 + \frac{\operatorname{Cov}(x_1, u)}{\operatorname{Var}(x_1)}$$

The term $\operatorname{Cov}(x_1, u)/\operatorname{Var}(x_1)$ represents the \textbf{asymptotic bias}. If $\operatorname{Cov}(x_1, u) > 0$, the estimator is asymptotically positively biased.
\subsection{Asymptotic Normality and Large-Sample Inference}

Consistency alone does not allow for hypothesis testing or confidence intervals. For inference, we need to know the sampling distribution of the estimators. Without MLR.6, the OLS estimators are generally not normally distributed in small samples. However, we can appeal to the \textbf{Central Limit Theorem (CLT)} to show that they are approximately normally distributed in large samples.

\begin{theorem}[Asymptotic Normality of OLS]
Under the Gauss-Markov assumptions (MLR.1 through MLR.5), as the sample size $n$ grows:
$$\frac{\hat{\beta}_j - \beta_j}{se(\hat{\beta}_j)} \xrightarrow{d} \text{Normal}(0, 1)$$
\end{theorem}

This theorem is powerful because it drops the strict normality assumption (MLR.6). It implies that $t$-statistics and $F$-statistics have approximate $t$ and $F$ distributions in large samples, regardless of the population distribution of $u$. Consequently, we can use the same testing procedures developed in this chapter, treating them as asymptotically valid approximations.

Note that this result still relies on the homoskedasticity assumption (MLR.5). If the error variance is not constant, the standard errors and test statistics derived here are invalid even in large samples. In such cases, robust standard error procedures become necessary.
\subsection{The Lagrange Multiplier Statistic}

For testing multiple exclusion restrictions in large samples, an alternative to the $F$-test is the \textbf{Lagrange Multiplier (LM) statistic}. The LM test only requires estimating the \textit{restricted model}, which can be computationally advantageous.

To test $H_0: \beta_{k-q+1} = \cdots = \beta_k = 0$ (excluding $q$ variables):
\begin{enumerate}
    \item Regress $y$ on the restricted set of independent variables and save the residuals $\hat{u}$.
    \item Regress $\hat{u}$ on \emph{all} independent variables (both included and excluded) and obtain the $R^2$ from this auxiliary regression, $R_{\text{aux}}^2$.
    \item Compute the test statistic:
\end{enumerate}
$$LM = n \cdot R_{\text{aux}}^2 \xrightarrow{d} \chi^2_q$$

Under $H_0$, $LM$ is distributed asymptotically as chi-squared with $q$ degrees of freedom. We reject $H_0$ if $LM$ exceeds the critical value from the $\chi^2_q$ distribution.

The intuition is that if the omitted variables truly have no effect, they should not be correlated with the residuals of the restricted model; thus $R_{\text{aux}}^2$ should be close to zero. A high $R_{\text{aux}}^2$ suggests the omitted variables do explain significant variation in the residuals, leading to rejection of the null.

% ===============================

% ===============================
\section{Robust Regression: Heteroskedasticity and Quantile Methods}
% ===============================

The preceding sections developed the core inferential machinery of regression under ideal conditions. We now relax those conditions in directions that are especially relevant for financial data: heteroskedastic disturbances, non-Gaussian tails, and the need to describe conditional distributions beyond the mean.
\subsection{Definition and Economic Motivation}

Heteroskedasticity occurs when the variance of the error term is not constant across observations. Formally:
$$\operatorname{Var}(u \mid \mathbf{x}) \neq \sigma^2$$

This violation is common in empirical applications: in financial data, returns on smaller investments exhibit higher volatility; in labor economics, wage variability increases with education; in firm-level data, profit dispersion widens with company size. These patterns reflect genuine economic heterogeneity---different subgroups in the data face different sources of uncertainty---rather than statistical aberrations.

While heteroskedasticity is prevalent in real-world datasets, its presence has critical implications for statistical inference. The OLS point estimates $\hat{\beta}_j$ remain \textbf{unbiased and consistent} under heteroskedasticity (because the derivation of OLS does not rely on the homoskedasticity assumption). However, the standard error formulas produced by conventional OLS regression are \textbf{biased}, invalidating $t$-statistics, confidence intervals, and hypothesis tests. In practice, heteroskedasticity more often produces underestimated standard errors, leading researchers to declare false statistical significance.
\subsection{Consequences for OLS Inference}

Under homoskedasticity, the variance of the OLS slope estimator is:
$$\operatorname{Var}(\hat{\beta}_j \mid \mathbf{x}) = \frac{\sigma^2}{SST_j(1 - R_j^2)}$$

When heteroskedasticity is present ($\operatorname{Var}(u \mid \mathbf{x}) = h(\mathbf{x})$ for some function $h$), this formula no longer applies. The practical consequence is severe: $t$-statistics computed using conventional standard errors do not follow a $t$-distribution under heteroskedasticity. Hypothesis tests conducted at a nominal 5\% significance level may reject the null at actual rates far different from 5\%.

Furthermore, OLS ceases to be BLUE. The Gauss-Markov Theorem, which guarantees that OLS is BLUE, explicitly requires homoskedasticity. When this assumption is violated, alternative estimators such as Weighted Least Squares can extract more information from the data.
\subsection{Testing for Heteroskedasticity: The White Test}

The White test provides a general diagnostic for heteroskedasticity. The procedure is:

\begin{enumerate}
\item Estimate the regression by OLS and compute residuals $\hat{u}_i$.
\item Regress squared residuals $\hat{u}_i^2$ on the original explanatory variables, their squares, and their pairwise interactions:
$$\hat{u}_i^2 = \gamma_0 + \gamma_1 x_1 + \cdots + \gamma_k x_k + \gamma_{k+1} x_1^2 + \gamma_{k+2} x_2^2 + \cdots + \gamma_m x_1 x_2 + v_i$$
\item Compute $R^2_{\text{aux}}$ from this auxiliary regression.
\item The test statistic is:
$$\text{WT} = n \cdot R^2_{\text{aux}} \xrightarrow{d} \chi^2_k$$
where $k$ is the number of regressors in the auxiliary regression (excluding the constant). Under the null hypothesis of homoskedasticity, $\text{WT}$ follows approximately a chi-squared distribution.
\end{enumerate}

If the test statistic exceeds the critical value, we reject homoskedasticity. The White test is robust because it makes no assumption about the specific functional form of heteroskedasticity; it simply tests whether squared residuals are systematically correlated with the explanatory variables and their nonlinear transformations.

Visual inspection of residual plots also provides useful diagnostics. A plot of standardized residuals $\hat{u}_i/\hat{\sigma}$ against fitted values $\hat{y}_i$ should show no systematic pattern under homoskedasticity. A widening or narrowing cone shape indicates heteroskedasticity correlated with the fitted value.

% Asset/code block omitted pending provenance acceptance.
\subsection{Heteroskedasticity-Robust Inference}

The standard solution in modern applied econometrics is to retain OLS point estimates but adjust standard errors to account for heteroskedasticity of unknown form. White's heteroskedasticity-robust variance estimator provides the ``sandwich'' variance-covariance matrix:
$$\widehat{\operatorname{Var}}(\hat{\boldsymbol{\beta}}) = (\mathbf{X}'\mathbf{X})^{-1} \left(\sum_{i=1}^n \hat{u}_i^2\, \mathbf{x}_i\mathbf{x}_i'\right) (\mathbf{X}'\mathbf{X})^{-1}$$

These ``Huber-White'' or ``sandwich'' standard errors are asymptotically valid under heteroskedasticity of any form. Under homoskedasticity, they remain valid, though typically slightly larger than conventional standard errors. This asymmetry makes robust standard errors the pragmatic default in applied work.

Once robust standard errors are computed, $t$-statistics and confidence intervals are constructed in the usual way:
$$t_j = \frac{\hat{\beta}_j}{se_{\text{robust}}(\hat{\beta}_j)}, \quad \text{CI}_j = \left[\hat{\beta}_j \pm 1.96 \cdot se_{\text{robust}}(\hat{\beta}_j)\right]$$

A coefficient may shift from statistically significant under conventional standard errors to insignificant under robust standard errors (or vice versa), reflecting a more honest appraisal of uncertainty in the presence of heteroskedasticity.
\subsection{Weighted Least Squares: An Alternative}

When the heteroskedasticity pattern is known or can be reliably estimated, Weighted Least Squares (WLS) offers improved efficiency. WLS downweights observations with high error variance and upweights observations with low error variance:
$$\min_{\boldsymbol{\beta}} \sum_{i=1}^n w_i (y_i - \mathbf{x}_i'\boldsymbol{\beta})^2, \quad w_i = 1/\hat{\sigma}_i^2$$

The WLS estimator is more efficient than OLS when the variance function is correctly specified. In practice, however, misspecification of the variance function can lead to biased or inefficient estimates. For this reason, heteroskedasticity-robust OLS has become the standard approach.

% ===============================

Traditional regression analysis, centered on OLS, provides a grand summary for the averages of the conditional distributions. While powerful, the conditional mean offers only a partial view of the relationship between variables. Just as the mean gives an incomplete picture of a single distribution, the regression curve gives a correspondingly incomplete picture for a set of distributions. \textbf{Quantile regression} offers a comprehensive strategy for completing this regression picture.

% ---- next approved source slice ----

\section{Financial Time Series: Returns, Dependence, and Linear Processes}
% ===============================

We now move fully into financial time series. The presentation begins with returns and stylized facts, then studies dependence through stationarity, correlation structures, white noise representations, and the first linear models used to describe predictable components of returns.
\subsection{Returns: Definitions and Properties}

To effectively analyze financial time series, it is essential to first introduce some fundamental financial definitions. A financial time series is a sequence of data points---such as asset prices, returns, or interest rates---recorded at regular time intervals. These series are central to understanding market behavior, risk, and forecasting in finance.

\begin{definition}[One-Period Simple Return]
Suppose we are examining an asset with price $P_t$ at time $t$. The \emph{simple gross return} over the period from $t-1$ to $t$ is:
$$1 + R_t = \frac{P_t}{P_{t-1}} \quad \text{or equivalently} \quad P_t = P_{t-1}(1 + R_t)$$
\end{definition}

\begin{definition}[Multiperiod Simple Return]
Over $k$ periods between $t-k$ and $t$:
$$1 + R_t[k] = \frac{P_t}{P_{t-k}} = \frac{P_t}{P_{t-1}} \times \frac{P_{t-1}}{P_{t-2}} \times \cdots \times \frac{P_{t-k+1}}{P_{t-k}} = \prod_{j=0}^{k-1}(1 + R_{t-j})$$
\end{definition}

To better analyze statistical properties, the \textbf{log return} (or continuously compounded return) is used in practice.

\begin{definition}[Continuously Compounded Return]
$$r_t = \ln(1 + R_t) = \ln\frac{P_t}{P_{t-1}} = p_t - p_{t-1}$$
where $p_t = \ln(P_t)$. Extending to multiperiod returns:
$$r_t[k] = \ln(1 + R_t[k]) = r_t + r_{t-1} + \cdots + r_{t-k+1}$$
The main advantage is \textbf{time additivity}: the total log return over multiple periods is simply the sum of the individual log returns, making statistical analysis more tractable.
\end{definition}
\subsection{Distributional Features of Returns}

Consider $N$ assets held for $T$ periods. For each asset $i$, let $r_{it}$ be the log return at time $t$. The log returns $\{r_{it}\}$ are often treated as i.i.d. as a baseline model, although this assumption has important limitations.

The most general model for the joint distribution of log returns is the conditional CDF:
$$F_r(r_{11}, \ldots, r_{NT}; Y; \theta)$$

where $Y$ is a vector of state variables and $\theta$ is a vector of parameters. Treating $Y$ as known, we partition the joint distribution as:
$$F(r_{i1}, \ldots, r_{iT}; \theta) = F(r_{i1}) \prod_{t=2}^T F(r_{it} \mid r_{i,t-1}, \ldots, r_{i1})$$

This highlights the time-based dependence of log returns. Under normality, the conditional distribution is:
$$f(r_1, \ldots, r_T; \theta) = f(r_1; \theta) \prod_{t=2}^T \frac{1}{\sqrt{2\pi\sigma_t}} \exp\!\left(\frac{-(r_t - \mu_t)^2}{2\sigma_t^2}\right)$$

Empirical evidence consistently shows that financial returns deviate from normality: they exhibit excess kurtosis (fat tails) and skewness. These \textbf{stylized facts} motivate the more sophisticated time series models discussed in subsequent chapters.

% ===============================
\subsection{Stationarity}

The foundation of time series analysis is stationarity. A time series $\{r_t\}$ is \textit{strictly stationary} if the joint distribution of $(r_{t_1}, \ldots, r_{t_k})$ is invariant under time shift. This is a very strong condition that is hard to verify empirically.

A weaker and more commonly used version is \textit{weak stationarity} (or covariance stationarity). A time series $\{r_t\}$ is weakly stationary if:
\begin{itemize}
    \item[(a)] $E(r_t) = \mu$, a constant independent of $t$
    \item[(b)] $\operatorname{Cov}(r_t, r_{t-\ell}) = \gamma_\ell$, depending only on the lag $\ell$, not on $t$
\end{itemize}

In practice, weak stationarity implies that a time plot shows values fluctuating with constant variation around a fixed level. The covariance $\gamma_\ell = \operatorname{Cov}(r_t, r_{t-\ell})$ is called the lag-$\ell$ \textbf{autocovariance}, satisfying $\gamma_0 = \operatorname{Var}(r_t)$ and $\gamma_{-\ell} = \gamma_\ell$.

It is common to assume that asset return series are weakly stationary.
\subsection{Correlation and Autocorrelation Functions}

\begin{definition}[Correlation Coefficient]
The correlation coefficient between two random variables $X$ and $Y$ is:
$$\rho_{x,y} = \frac{\operatorname{Cov}(X,Y)}{\sqrt{\operatorname{Var}(X)\operatorname{Var}(Y)}} = \frac{E[(X-\mu_x)(Y-\mu_y)]}{\sqrt{E(X-\mu_x)^2 \cdot E(Y-\mu_y)^2}}$$
\end{definition}

This coefficient measures the strength of linear dependence, with $-1 \leq \rho_{x,y} \leq 1$. The two variables are uncorrelated if $\rho_{x,y} = 0$. Furthermore, if both $X$ and $Y$ are normal random variables, then $\rho_{x,y} = 0$ if and only if $X$ and $Y$ are independent. The sample counterpart is:
$$\hat{\rho}_{x,y} = \frac{\sum_{t=1}^T (x_t - \bar{x})(y_t - \bar{y})}{\sqrt{\sum_{t=1}^T (x_t - \bar{x})^2 \sum_{t=1}^T (y_t - \bar{y})^2}}$$

\begin{definition}[Autocorrelation Function (ACF)]
For a weakly stationary return series $r_t$, the lag-$\ell$ autocorrelation is:
$$\rho_\ell = \frac{\operatorname{Cov}(r_t, r_{t-\ell})}{\operatorname{Var}(r_t)} = \frac{\gamma_\ell}{\gamma_0}$$
\end{definition}

From the definition, $\rho_0 = 1$, $\rho_\ell = \rho_{-\ell}$, and $-1 \leq \rho_\ell \leq 1$. A weakly stationary series $r_t$ is not serially correlated if and only if $\rho_\ell = 0$ for all $\ell > 0$.

For a given sample $\{r_t\}_{t=1}^T$, the lag-$\ell$ sample autocorrelation is:
$$\hat{\rho}_\ell = \frac{\sum_{t=\ell+1}^T (r_t - \bar{r})(r_{t-\ell} - \bar{r})}{\sum_{t=1}^T (r_t - \bar{r})^2}, \quad 0 \leq \ell < T - 1$$

If $\{r_t\}$ is an i.i.d. sequence with $E(r_t^2) < \infty$, then $\hat{\rho}_\ell$ is asymptotically normal with mean zero and variance $1/T$.
\subsection{White Noise and Linear Time Series}

A time series $r_t$ is called a \textbf{white noise} if $\{r_t\}$ is a sequence of independent and identically distributed random variables with finite mean and variance. If $r_t$ is normally distributed with mean zero and variance $\sigma^2$, the series is a \textbf{Gaussian white noise}. For a white noise series, all ACFs are zero.

A time series $\{r_t\}$ is said to be \textbf{linear} if it can be written as:
$$r_t = \mu + \sum_{i=0}^\infty \psi_i a_{t-i}$$

where $\mu$ is the mean, $\psi_0 = 1$, and $\{a_t\}$ is a white noise series (the innovation or shock at time $t$).

Under stationarity:
$$E(r_t) = \mu, \quad \operatorname{Var}(r_t) = \sigma_a^2 \sum_{i=0}^\infty \psi_i^2, \quad \rho_\ell = \frac{\sum_{i=0}^\infty \psi_i \psi_{i+\ell}}{1 + \sum_{i=1}^\infty \psi_i^2}$$

For a weakly stationary time series, $\psi_i \to 0$ as $i \to \infty$, meaning that the linear dependence of $r_t$ on the remote past diminishes for large $\ell$.

% ===============================
\subsection{Autoregressive Models}
\subsubsection{The AR(1) Model}

For a series with significant lag-1 autocorrelation, the AR(1) model is:
$$r_t = \phi_0 + \phi_1 r_{t-1} + a_t$$

where $\{a_t\}$ is Gaussian white noise with zero mean and variance $\sigma_a^2$.

For weak stationarity, we need $|\phi_1| < 1$, which ensures that past shocks decay geometrically over time. Under this condition:
\begin{itemize}
    \item Unconditional mean: $\mu = \phi_0/(1 - \phi_1)$
    \item Unconditional variance: $\operatorname{Var}(r_t) = \sigma_a^2/(1 - \phi_1^2)$
    \item Autocovariances: $\gamma_\ell = \phi_1^\ell \gamma_0$
    \item ACF: $\rho_\ell = \phi_1^\ell$ (exponential decay)
\end{itemize}

The ACF decays monotonically for $\phi_1 > 0$, and alternates in sign for $\phi_1 < 0$.
\subsubsection{The AR(2) Model}

An AR(2) model:
$$r_t = \phi_0 + \phi_1 r_{t-1} + \phi_2 r_{t-2} + a_t$$

Stationarity requires all roots of $1 - \phi_1 z - \phi_2 z^2 = 0$ to lie outside the unit circle. The unconditional mean is $\mu = \phi_0/(1 - \phi_1 - \phi_2)$, and the ACF satisfies:
$$\rho_\ell = \phi_1 \rho_{\ell-1} + \phi_2 \rho_{\ell-2} \quad (\ell \geq 2), \qquad \rho_1 = \frac{\phi_1}{1-\phi_2}$$

Depending on the nature of the characteristic roots, the ACF shows either monotone exponential decay (real roots) or a damped sinusoidal pattern (complex conjugate roots).

\subsection[The AR(p) Model]{The AR(\texorpdfstring{$p$}{p}) Model}

The general AR($p$) model:
$$r_t = \phi_0 + \phi_1 r_{t-1} + \cdots + \phi_p r_{t-p} + a_t$$

requires all roots of the characteristic equation $1 - \phi_1 x - \cdots - \phi_p x^p = 0$ to be greater than one in modulus. The ACF satisfies:
$$(1 - \phi_1 B - \phi_2 B^2 - \cdots - \phi_p B^p)\rho_\ell = 0 \quad \text{for } \ell > 0$$

\textbf{Identifying AR Models: PACF.} The \textbf{partial autocorrelation function (PACF)} is the key tool for AR order identification. The lag-$j$ PACF $\hat{\phi}_{j,j}$ measures the added contribution of $r_{t-j}$ to $r_t$ over the AR$(j-1)$ model. For a stationary AR($p$) model, the sample PACF \textbf{cuts off at lag $p$}: $\hat{\phi}_{p,p}$ converges to $\phi_p$, while $\hat{\phi}_{\ell,\ell} \to 0$ for all $\ell > p$. In practice, approximate 95\% confidence limits of $\pm 2/\sqrt{T}$ are used to determine which lags are statistically significant.
\subsection{Moving-Average Models}

Unlike AR models, MA models always exhibit weak stationarity because they are finite linear combinations of a white noise sequence. An MA($q$) model is:
$$r_t = c_0 + a_t - \theta_1 a_{t-1} - \cdots - \theta_q a_{t-q}$$

\textbf{Properties:}
\begin{itemize}
    \item Mean: $E(r_t) = c_0$ (time-invariant)
    \item Variance: $\operatorname{Var}(r_t) = (1 + \theta_1^2 + \cdots + \theta_q^2)\sigma_a^2$
    \item The ACF \textbf{cuts off at lag $q$}: $\rho_\ell = 0$ for $\ell > q$
\end{itemize}

For invertibility, all roots of $1 - \theta_1 z - \cdots - \theta_q z^q = 0$ must lie outside the unit circle, so that the model can be expressed as an infinite AR process.

\textbf{Identifying MA Order.} The ACF is the key diagnostic: if the sample ACF shows $\rho_q \neq 0$ but $\rho_\ell \approx 0$ for $\ell > q$, the series follows an MA($q$) model. This is the dual of the PACF identification rule for AR models: AR models use PACF for order selection, MA models use ACF.
\subsection{Order Identification for AR and MA Models}
\subsubsection{Identifying AR Models}
In application, the order $p$ of an AR time series is unknown and must be specified empirically. This is referred to as the order determination of AR models, and it has been extensively studied in the time series literature.
In general, the partial autocorrelation function (PACF) is exploited for this purpose.\\
The PACF of a stationary time series is a useful tool for determining the order $p$ of an AR model. A simple yet effective way to introduce PACF is to consider AR models in consecutive orders. These models can be estimated by the least squares method as multiple linear regression. \\
The lag-$j$ sample PACF $\hat{\phi}_{j,j}$ measures the added contribution of $r_{t-j}$ to $r_t$ over the AR$(j-1)$ model. For a stationary Gaussian AR$(p)$ model, the sample PACF has the following properties:
\begin{itemize}
    \item $\hat{\phi}_{p,p}$ converges to $\phi_p$ as the sample size $T$ goes to infinity.
    \item $\hat{\phi}_{\ell,\ell}$ converges to zero for all $\ell > p$.
    \item The asymptotic variance of $\hat{\phi}_{\ell,\ell}$ is $1/T$ for $\ell > p$.
\end{itemize}

These properties imply that the sample PACF cuts off at lag $p$. In practice, we use this property to identify the order of an AR$(p)$ model: the lag-$p$ sample PACF should be significantly different from zero, whereas $\hat{\phi}_{j,j}$ should be close to zero for all $j > p$. Typically, we use approximate $95\%$ confidence limits of $\pm(2/\sqrt{T})$ to determine which lags are statistically significant.
\subsubsection{Identifying MA Order}
The ACF is useful in identifying the order of an MA model. If the sample ACF shows $\rho_q \neq 0$ but $\rho_\ell \approx 0$ for $\ell > q$, then the series follows an MA$(q)$ model.

Unlike AR models, which use the PACF for order selection, MA models use the sample ACF to directly identify the number of significant moving-average terms. The ACF of a stationary MA$(q)$ model cuts off at lag $q$: we observe significant autocorrelations up to lag $q$ and negligible autocorrelations beyond lag $q$.\\
In practice, approximate $95\%$ confidence limits for the sample ACF are $\pm(2/\sqrt{T})$. Lags with sample ACF outside these limits are considered statistically significant. The order $q$ is determined as the lag at which the ACF becomes negligible within the confidence bounds and remains negligible for subsequent lags.\\
Note that, unlike the PACF which provides information on the nonzero AR lags of the model, the sample ACF provides direct information on the nonzero MA lags. This fundamental difference reflects the dual nature of AR and MA representations: AR models are identified by the pattern of the PACF, while MA models are identified by the pattern of the ACF.

% ===============================
\section{ARMA, Nonstationarity, ARIMA, and Dynamic Correlation Models}
% ===============================

Having introduced the building blocks of linear dependence, we can now study models that combine autoregressive and moving-average components, deal with nonstationary price dynamics through differencing, and finally extend the analysis to time-varying dependence across assets through rolling and DCC correlations.
\subsection{ARMA Models}

In practice, a model combining both AR and MA structures may be more appropriate. An ARMA($p$,$q$) model:
$$r_t = \phi_0 + \phi_1 r_{t-1} + \cdots + \phi_p r_{t-p} + a_t - \theta_1 a_{t-1} - \cdots - \theta_q a_{t-q}$$

is stationary if all AR roots lie outside the unit circle, and invertible if all MA roots lie outside the unit circle. The unconditional mean under stationarity is $E(r_t) = \phi_0/(1 - \phi_1 - \cdots - \phi_p)$.

A key property: \textbf{both ACF and PACF tail off to zero}, without sharp cutoff. This signature characteristic distinguishes ARMA from pure AR (PACF cuts off) and pure MA (ACF cuts off) models.

\textbf{Model Identification via Information Criteria.} Since visual inspection is often ambiguous for ARMA models, a systematic approach uses:
$$\text{AIC}(p,q) = -2\log(\hat{L}) + 2(p + q + 1), \qquad \text{BIC}(p,q) = -2\log(\hat{L}) + (p + q + 1)\log(T)$$

The preferred model minimizes the criterion. A practical strategy fits a range of ARMA($p$,$q$) models, selects the one with minimum AIC or BIC, and verifies adequacy by checking that the residuals behave as white noise. BIC tends to select more parsimonious models than AIC.

% ===============================
\subsubsection{The ARMA(1,1) Model}
\textbf{Stationarity and Invertibility.}
For an ARMA(1,1) model to be weakly stationary, the AR component must satisfy $|\phi_1| < 1$. For the model to be invertible, the MA component must satisfy $|\theta_1| < 1$. Under these conditions, the process admits both a stationary AR representation with infinite order and an invertible MA representation with infinite order.

The unconditional mean of an ARMA(1,1) model is
\[
\mathbb{E}(r_t) = \frac{\phi_0}{1 - \phi_1}.
\]
The unconditional variance is
\[
\mathrm{Var}(r_t) = \frac{(1 - 2\phi_1\theta_1 + \theta_1^2)}{(1 - \phi_1^2)}\sigma_a^2.
\]

\textbf{Autocorrelation Function.}
Unlike pure AR or MA models, the ACF of an ARMA(1,1) model does not exhibit a simple cutoff pattern. Instead, the ACF decays exponentially after the initial lags. Specifically, for lags $\ell \geq 2$,
\[
\rho_\ell = \phi_1 \rho_{\ell-1},
\]
which shows that the ACF follows an AR(1)-like recursive structure for $\ell \geq 2$. The lag-1 ACF is
\[
\rho_1 = \frac{\phi_1(1 - \phi_1\theta_1) + (\theta_1 - \phi_1)}{1 + \theta_1^2 - 2\phi_1\theta_1}.
\]
The PACF of an ARMA(1,1) model also exhibits an exponential decay pattern similar to that of an MA(1) model.

\subsection[General ARMA(p,q) Models]
{General ARMA(\texorpdfstring{$p$}{p},\texorpdfstring{$q$}{q}) Models}
The general ARMA$(p,q)$ model combines $p$ autoregressive terms and $q$ moving-average terms. The model is stationary if all roots of the characteristic AR polynomial $1 - \phi_1 z - \cdots - \phi_p z^p = 0$ lie outside the unit circle. The model is invertible if all roots of the MA characteristic polynomial $1 - \theta_1 z - \cdots - \theta_q z^q = 0$ lie outside the unit circle.

Under stationarity and invertibility conditions, the unconditional mean is
\[
\mathbb{E}(r_t) = \frac{\phi_0}{1 - \phi_1 - \cdots - \phi_p}.
\]

A key property of ARMA$(p,q)$ models is that:
\begin{itemize}
    \item The ACF decays exponentially with a pattern determined by the AR characteristic roots. It does not exhibit the sharp cutoff of pure AR or MA models.
    \item The PACF also decays exponentially, without sharp cutoff.
    \item Both ACF and PACF tail off to zero, indicating a mixed structure.
\end{itemize}

This property of both ACF and PACF tailing off is the signature characteristic of ARMA models and distinguishes them from pure AR and MA models.
\subsubsection{Identifying ARMA Models}
Identifying the orders $p$ and $q$ of an ARMA model is more challenging than identifying pure AR or MA models because both the ACF and PACF exhibit exponential decay without sharp cutoff. Several approaches are available for model identification.

\textbf{Visual Inspection of ACF and PACF.}
The sample ACF and PACF provide visual clues for ARMA identification. If both the sample ACF and PACF decay exponentially without showing a clear cutoff, an ARMA model is suggested. A plot showing ACF tailing off gradually and PACF also tailing off gradually is indicative of an ARMA structure.

\textbf{Information Criteria.}
A more systematic approach to identify ARMA orders is to use information criteria such as the Akaike Information Criterion (AIC) or the Bayesian Information Criterion (BIC):
\[
\text{AIC}(p,q) = -2\log(\text{likelihood}) + 2(p + q + 1),
\]
\[
\text{BIC}(p,q) = -2\log(\text{likelihood}) + (p + q + 1)\log(T),
\]
where $T$ is the sample size. The preferred ARMA$(p,q)$ model is the one that minimizes the information criterion over a range of candidate values.

\textbf{Practical Strategy.}
In practice, a sequential approach is often employed:
\begin{enumerate}
    \item Fit a range of ARMA$(p,q)$ models with $p$ and $q$ typically ranging from 0 to 3 or 4.
    \item Compute AIC or BIC for each candidate model.
    \item Select the model with the minimum information criterion value.
    \item Verify the adequacy of the selected model by examining the residuals to ensure they behave as white noise.
\end{enumerate}

This approach balances model fit with parsimony, avoiding both underfitting and overfitting.

So far we have focused on return series that are stationary. In some studies, interest rates, foreign exchange rates, or price series of assets are of interest. These series tend to be nonstationary. For a price series, the nonstationarity is mainly due to the fact that there is no fixed level for the price. Such a nonstationary series is called a \textit{unit-root nonstationary} time series.
\subsection{Random Walk and Drift}

A time series $\{p_t\}$ is a random walk if:
$$p_t = p_{t-1} + a_t$$

where $p_0$ is the starting value and $\{a_t\}$ is white noise. Treating this as a special AR(1) model, the coefficient of $p_{t-1}$ is unity, which violates the stationarity condition $|\phi_1| < 1$. A random-walk series is therefore unit-root nonstationary.

\textbf{Forecasting Properties.} The $\ell$-step ahead forecast from origin $h$ is:
$$\hat{p}_h(\ell) = E(p_{h+\ell} \mid p_h, p_{h-1}, \ldots) = p_h \quad \forall\, \ell > 0$$

Point forecasts are simply the current value. The forecast error variance $\operatorname{Var}[e_h(\ell)] = \ell\sigma_a^2$ diverges to infinity as $\ell \to \infty$.

\textbf{Long Memory and Permanent Shocks.} The MA representation is:
$$p_t = \sum_{i=0}^{t-1} a_{t-i}$$

All coefficients $\psi_i = 1$ for all $i$. The impact of any past shock $a_{t-i}$ on $p_t$ does not decay over time: shocks have a \textbf{permanent effect}. This contrasts with stationary series, where the influence of past shocks decays geometrically.

The random walk with drift:
$$p_t = \mu + p_{t-1} + a_t$$

includes the constant $\mu$, which represents the time trend of the log price. A positive drift ($\mu > 0$) implies that the log price eventually goes to infinity.
\subsection{Trend-Stationary Time Series}

A closely related model exhibiting linear trend is the trend-stationary time series:
$$p_t = \beta_0 + \beta_1 t + r_t$$

where $r_t$ is a stationary AR($p$) series. Despite similar appearance to a random walk with drift, there is a fundamental difference:
\begin{itemize}
    \item \textbf{Random walk with drift}: $\operatorname{Var}(p_t) = t\sigma_a^2$, which grows without bound.
    \item \textbf{Trend-stationary model}: $\operatorname{Var}(p_t) = \operatorname{Var}(r_t)$, which is finite and time-invariant.
\end{itemize}

The trend-stationary series can be transformed into a stationary one by removing the time trend via linear regression analysis.
\subsection{ARIMA Models and Differencing}

An ARIMA model extends ARMA by allowing the AR polynomial to have 1 as a characteristic root. A time series $y_t$ is an ARIMA($p$, 1, $q$) process if the change series:
$$c_t = y_t - y_{t-1} = (1 - B)y_t$$

follows a stationary and invertible ARMA($p$,$q$) model. In finance, price series are commonly believed to be nonstationary, but the log return series $r_t = \ln(p_t) - \ln(p_{t-1})$ is stationary. The log price is therefore unit-root nonstationary and can be treated as an ARIMA process.

Like a random walk, an ARIMA model has strong memory because the $\psi_i$ coefficients in its MA representation do not decay to zero, implying that past shocks have a permanent effect on the series.
\subsection{The Transition to GARCH and the Limitations of Linear Models}

ARMA and ARIMA models are designed to model the conditional mean $E(r_t \mid r_{t-1}, r_{t-2}, \ldots)$ of financial return series. While these frameworks capture the linear dependence structure in returns, they suffer from a fundamental limitation: they assume that the conditional variance is \textbf{constant} over time:
$$\operatorname{Var}(r_t \mid r_{t-1}, \ldots) = \sigma_a^2$$

In practice, financial returns exhibit \textbf{volatility clustering}: large price changes tend to be followed by large changes, and small changes tend to be followed by small changes. This stylized fact cannot be captured by linear models. More broadly, linear models are limited in the following ways:

\begin{itemize}
    \item \textbf{Constant conditional variance:} They cannot model time-varying volatility, essential for risk management and derivative pricing.
    \item \textbf{Symmetric response:} Linear models treat positive and negative shocks of equal magnitude identically, while financial returns often exhibit asymmetric responses (the leverage effect).
    \item \textbf{No nonlinear dependence:} The autocorrelation structure of raw returns may appear weak, but squared returns $r_t^2$ often exhibit strong autocorrelation---evidence of nonlinear dependence that ARMA cannot capture.
    \item \textbf{Gaussian tails:} Linear models with Gaussian errors underestimate the probability of extreme events.
\end{itemize}

The solution is to apply the same autoregressive logic that ARMA uses for conditional means to modeling conditional variances instead. A GARCH($p$,$q$) model:
$$\sigma_t^2 = \alpha_0 + \sum_{i=1}^{q} \alpha_i a_{t-i}^2 + \sum_{j=1}^{p} \beta_j \sigma_{t-j}^2$$

expresses the conditional variance as a function of past squared shocks and past variances. The architectural similarity between ARMA and GARCH is striking: both embody the same autoregressive principle applied to different moments of the distribution. ARMA and ARIMA models are the essential foundation upon which modern volatility modeling is built. The ARCH and GARCH family---to be treated in detail in forthcoming articles---extends this framework to accommodate all the stylized facts of financial time series.

% ===============================

The preceding sections focused on univariate time series---the evolution of a single asset over time. In practice, understanding how the relationships \emph{between} assets evolve over time is equally important. A portfolio manager needs to know not only the volatility of individual assets but also how their correlations shift during market stress. This chapter introduces two key tools for studying dynamic dependence: rolling correlation and dynamic conditional correlation.
\subsection{Rolling Correlation}

The simplest approach to time-varying correlation is the \textbf{rolling (window) correlation}. The idea is to compute the standard Pearson correlation between two return series over a rolling window of fixed length $w$, updating the estimate at each time step:
$$\hat{\rho}_t(r_x, r_y) = \frac{\sum_{s=t-w+1}^{t}(r_{x,s} - \bar{r}_x)(r_{y,s} - \bar{r}_y)}{\sqrt{\sum_{s=t-w+1}^{t}(r_{x,s} - \bar{r}_x)^2 \sum_{s=t-w+1}^{t}(r_{y,s} - \bar{r}_y)^2}}$$

where $\bar{r}_x$ and $\bar{r}_y$ are the sample means within the window $[t-w+1, t]$.

Rolling correlation captures how the linear relationship between two assets evolves through time, revealing dynamics such as flight-to-quality (when equity-bond correlations turn negative during market crises) or contagion effects (when cross-market correlations spike during stress periods).

\begin{remark}
The choice of window length $w$ involves a bias-variance tradeoff. A short window is more responsive to recent changes but noisier. A long window is smoother but slow to detect structural breaks. Practitioners typically experiment with windows ranging from 30 to 252 trading days depending on the data frequency and the phenomenon of interest.
\end{remark}

The main limitation of rolling correlation is that it weights all observations within the window equally, regardless of recency. Furthermore, it has no theoretical structure---it is a purely descriptive tool without an underlying model for the correlation dynamics.
\subsection{Dynamic Conditional Correlation (DCC)}

The \textbf{Dynamic Conditional Correlation (DCC)} model, introduced by Engle (2002), provides a theoretically grounded alternative. DCC extends the GARCH framework to the multivariate setting, allowing the correlation matrix to evolve over time.

The DCC model decomposes the conditional covariance matrix $\mathbf{H}_t$ as:
$$\mathbf{H}_t = \mathbf{D}_t \mathbf{R}_t \mathbf{D}_t$$

where:
\begin{itemize}
    \item $\mathbf{D}_t = \operatorname{diag}(\sqrt{h_{11,t}}, \ldots, \sqrt{h_{NN,t}})$ is a diagonal matrix of conditional standard deviations, with each $h_{ii,t}$ following its own univariate GARCH process
    \item $\mathbf{R}_t$ is the time-varying correlation matrix
\end{itemize}

The correlation matrix $\mathbf{R}_t$ is obtained from a proxy process $\mathbf{Q}_t$:
$$\mathbf{Q}_t = (1 - \alpha - \beta)\bar{\mathbf{Q}} + \alpha \boldsymbol{\varepsilon}_{t-1}\boldsymbol{\varepsilon}_{t-1}' + \beta \mathbf{Q}_{t-1}$$

where $\bar{\mathbf{Q}}$ is the unconditional covariance matrix of the standardized residuals $\boldsymbol{\varepsilon}_t = \mathbf{D}_t^{-1}\mathbf{r}_t$, and $\alpha, \beta \geq 0$ with $\alpha + \beta < 1$. The conditional correlation matrix is then:
$$\mathbf{R}_t = \mathbf{Q}_t^{*-1} \mathbf{Q}_t \mathbf{Q}_t^{*-1}$$

where $\mathbf{Q}_t^* = \operatorname{diag}(\sqrt{q_{11,t}}, \ldots, \sqrt{q_{NN,t}})$ rescales $\mathbf{Q}_t$ so that each diagonal element of $\mathbf{R}_t$ equals one.

The DCC model is estimated in two steps. First, univariate GARCH models are fitted to each asset to obtain the conditional standard deviations and the standardized residuals. Second, the DCC parameters $\alpha$ and $\beta$ are estimated by maximizing the correlation part of the log-likelihood.

\begin{remark}
The parameter $\alpha$ measures the sensitivity of current correlations to recent innovations (news impact on correlation), while $\beta$ measures the persistence of past correlations. The stationarity condition $\alpha + \beta < 1$ ensures that correlations mean-revert to the unconditional level $\bar{\mathbf{Q}}$ over time.
\end{remark}

The DCC model offers several advantages over rolling correlation: it provides a maximum-likelihood framework, it respects the positive-definiteness constraint on correlation matrices, and it accommodates asymmetric responses through extensions such as the Asymmetric DCC (ADCC). For analyzing the dynamic relationship between gold and other financial assets, the DCC model provides a principled and flexible framework.


% ===============================
